% 第7章 金属中的磁性（Magnetism in metals）
\chapter{金属中的磁性}

本章讨论金属的磁性质。前几章我们集中于有相互作用但局域的磁矩。金属中的传导电子是
离域的，可以在样品中自由游走——它们称为巡游电子（itinerant electrons）。有些
情形金属中的磁矩与传导电子相联系；另一些情形磁矩保持局域。两种情形下都可以出现
顺磁与抗磁行为；一定条件下铁磁也可能。本章的大部分讨论将围绕自由电子模型展开
（下一节引入）。自由电子模型对多数真实情形是粗糙的近似，但便于处理，能使讨论
走得很远。后续各节推导电子气磁性质：包括泡利顺磁性、朗道抗磁性、RKKY 相互作用
的起源、电子气的不稳定性（如自旋密度波的形成），以及局域磁矩与电子气相互作用
时发生的近藤效应。

\section{自由电子模型}

通过回顾自由电子模型开始巡游电子磁性的讨论。此模型忽略晶格的周期势，电子填充
直到费米波矢 $k_{\mathrm{F}}$ 的各个态。k 空间中相邻点相距 $2\pi/L$
（见图 7.1(a)），$V=L^{3}$ 是样品体积，故 k 与 $k+\mathrm{d}k$ 之间的态数等于
$4\pi k^{2}\mathrm{d}k$（半径 k、宽 dk 的球壳体积，见图 7.1(b)）除以 k 空间中
一个点占据的体积 $(2\pi/L)^{3}$。每个态被双占据——一个自旋向上电子与一个自旋
向下电子——故再有因子 2。于是态密度 $g(k)\mathrm{d}k$ 可写为
\begin{equation*}
g(k)\,\mathrm{d}k = \frac{2}{(2\pi/L)^{3}}\times4\pi k^{2}\,\mathrm{d}k,\tag{7.1}
\end{equation*}
其中因子 2 计入电子的两个自旋态。于是
\begin{equation*}
g(k)\,\mathrm{d}k = \frac{Vk^{2}\,\mathrm{d}k}{\pi^{2}}.\tag{7.2}
\end{equation*}
若材料有 N 个电子，绝对零度（T=0）下它们将填满直到最大波矢 $k_{\mathrm{F}}$ 的态：
\begin{equation*}
N = \int_{0}^{k_{\mathrm{F}}}g(k)\,\mathrm{d}k = \frac{Vk_{F}^{3}}{3\pi^{2}},\tag{7.3}
\end{equation*}
故
\begin{equation*}
k_{\mathrm{F}}^{3} = 3\pi^{2}n\tag{7.4}
\end{equation*}
$n=N/V$ 是单位体积电子数。费米能 $E_{\mathrm{F}}$ 定义为
\begin{equation*}
E_{\mathrm{F}} = \frac{\hbar^{2}k_{\mathrm{F}}^{2}}{2m_{\mathrm{e}}}.\tag{7.5}
\end{equation*}
态密度作为能量（$E\propto k^{2}$）的函数正比于 $E^{1/2}$：
\begin{equation*}
g(E) = \frac{\mathrm{d}n}{\mathrm{d}E} \propto E^{1/2}\tag{7.6}
\end{equation*}
于是
\begin{equation*}
n = \int_{0}^{E_{\mathrm{F}}}g(E)\,\mathrm{d}E \propto E_{\mathrm{F}}^{3/2}\tag{7.7}
\end{equation*}
从而
\begin{equation*}
\frac{\mathrm{d}n}{n} = \frac{3}{2}\frac{\mathrm{d}E_{\mathrm{F}}}{E_{\mathrm{F}}}.\tag{7.8}
\end{equation*}
费米能处的态密度为\fn{1}{式 7.9 也可以直接用式 7.2、7.4、7.5 导出。}
\begin{equation*}
g(E_{\mathrm{F}}) = \left(\frac{\mathrm{d}n}{\mathrm{d}E}\right)_{E=E_{\mathrm{F}}} = \frac{3}{2}\frac{n}{E_{\mathrm{F}}}.\tag{7.9}
\end{equation*}
顺带指出：联立式 7.4、7.5 与 7.9 可得 $g(E_{\mathrm{F}})$ 的另一个有用表达式
\begin{equation*}
g(E_{\mathrm{F}}) = \frac{m_{\mathrm{e}}k_{\mathrm{F}}}{(\pi\hbar)^{2}},\tag{7.10}
\end{equation*}
它表明 $g(E_{\mathrm{F}})\propto m_{\mathrm{e}}$。许多重要性质依赖费米能处的态密度，
因此知道它正比于电子质量是有用的。在许多感兴趣的体系中，由于能带结构或相互作用的
效应，电子质量相对其自由空间值被增强。

自由电子模型忽略晶格周期势。若把它作为微扰计入（近自由电子模型），结果是：除了
电子波矢接近倒格矢之处，它几乎没有影响；在这些 k 空间点，色散关系中出现能隙
（见图 7.2）。

\begin{figure}
\centering
\includegraphics[width=0.4\textwidth]{fig7_01.pdf}
\caption{（a）电子态在 k 空间相距 $2\pi/L$；每个态可双占据、占据体积 $(2\pi/L)^{3}$。
（b）态密度的计算：考虑 k 空间中波矢 k 与 $k+\mathrm{d}k$ 的态之间的体积，即
$4\pi k^{2}\mathrm{d}k$。}
\end{figure}

\begin{figure}
\centering
\includegraphics[width=0.45\textwidth]{fig7_02.pdf}
\caption{布里渊区边界处的能隙。}
\end{figure}

迄今一切都按 T=0 处理。$T>0$ 时，态密度 $g(E)$ 不变，但每个态的占据由费米函数
$f(E)$ 决定：
\begin{equation*}
f(E) = \frac{1}{\mathrm{e}^{(E-\mu)/k_{\mathrm{B}}T}+1},\tag{7.11}
\end{equation*}
$\mu$ 是依赖温度的化学势。此函数绘于图 7.3。T=0 时 $f(E)$ 是阶跃函数：$E<\mu$
取 1、$E>\mu$ 取 0；温度升高时台阶被抹平。当费米函数接近阶跃函数——多数金属在熔点
以下几乎所有温度的通常情形——称电子处于简并极限（degenerate limit）。费米函数
是泡利不相容原理的结果：每个电子必须有一组唯一的量子数，没有两个电子能处于同一
态。$(E-\mu)\gg k_{\mathrm{B}}T$ 时费米函数趋于麦克斯韦--玻尔兹曼形式
$\mathrm{e}^{-(E-\mu)/k_{\mathrm{B}}T}$，称为非简并极限。

费米能是 T=0 时最高被占据能级的能量，由方程
\begin{equation*}
\int_{0}^{E_{\mathrm{F}}}f(E)g(E)\,\mathrm{d}E = n.\tag{7.12}
\end{equation*}
确定。函数 $f(E)g(E)$ 示于图 7.3(c)。T=0 时容易得到：费米能精确等于化学势
$E_{\mathrm{F}}=\mu$。$T>0$ 时经冗长的计算得
\begin{equation*}
\mu = E_{\mathrm{F}}\left[1 - \frac{\pi^{2}}{12}\left(\frac{k_{\mathrm{B}}T}{E_{\mathrm{F}}}\right)^{2} + O\left(\left(\frac{k_{\mathrm{B}}T}{E_{\mathrm{F}}}\right)^{4}\right)\right]\tag{7.13}
\end{equation*}
但这意味着：对典型金属，即使在室温，把 $E_{\mathrm{F}}$ 与 $\mu$ 等同起来也精确到
0.01\%——尽管值得记住这两个量并不相同。费米面是 k 空间中能量等于化学势的点集。
若化学势位于能隙内，材料是半导体或绝缘体，没有费米面。因此，金属就是有费米面的
材料。

\begin{figure}
\centering
\includegraphics[width=0.32\textwidth]{fig7_03.pdf}
\caption{（a）式 7.11 定义的费米函数 $f(E)$：粗线为 $T=0$；温度升高时台阶被抹平。
所标温度为 $T=0$、$T=0.01\mu/k_{\mathrm{B}}$、$T=0.05\mu/k_{\mathrm{B}}$ 与
$T=0.1\mu/k_{\mathrm{B}}$。（b）自由电子气的态密度 $g(E)$ 正比于 $E^{1/2}$。
（c）与（a）相同温度下的 $f(E)g(E)$。}
\end{figure}

\section{泡利顺磁性}

金属中每个 k 态因电子有两个可能的自旋态而可双占据，因此金属中每个电子要么自旋向上、要么自旋向下。
加磁场时，电子的能量随其自旋而升降。这给出电子气的顺磁磁化率，称为泡利顺磁性
（Pauli paramagnetism）。

\subsection{初等推导}

起初我们忽略轨道贡献、取 g=2，并忽略有限温度造成的费米面模糊。如图 7.4 所示，外加
磁场中电子能带按自旋分裂成两个子带，间隔 $g\mu_{\mathrm{B}}B=2\mu_{\mathrm{B}}B$。
设 $g\mu_{\mathrm{B}}B$ 是很小的能量，能带分裂非常小。自旋向上的额外电子数密度为
$n_{\uparrow}=\frac{1}{2}g(E_{\mathrm{F}})\mu_{\mathrm{B}}B$；自旋向下的亏损
数密度亦然：$n_{\downarrow}=\frac{1}{2}g(E_{\mathrm{F}})\mu_{\mathrm{B}}B$。于是
磁化强度为
\begin{equation*}
M = \mu_{\mathrm{B}}(n_{\uparrow}-n_{\downarrow}) = g(E_{\mathrm{F}})\mu_{\mathrm{B}}^{2}B\tag{7.14}
\end{equation*}
磁化率 $\chi_{\mathrm{P}}$（下标 P 表示泡利磁化率）为
\begin{align*}
\chi_{\mathrm{P}} = \frac{M}{H}\approx\frac{\mu_0M}{B} &= \mu_0\mu_{\mathrm{B}}^{2}g(E_{\mathrm{F}})\tag{7.15}\\
&= \frac{3n\mu_0\mu_{\mathrm{B}}^{2}}{2E_{\mathrm{F}}}\tag{7.16}
\end{align*}
最后一步用式 7.9。由于 $\chi_{\mathrm{P}}\ll1$，写 $\chi_{\mathrm{P}}\approx\mu_0M/B$
是合理的（见 1.1.4 节）。

\begin{figure}
\centering
\includegraphics[width=0.45\textwidth]{fig7_04.pdf}
\caption{态密度图显示磁场 B 中能带的分裂。分裂被大大夸大。}
\end{figure}

我们的泡利顺磁性表达式不依赖温度——诚然，这是因为我们一开始就忽略了有限温度的
费米面模糊。不过若计入温度，它只是一个相当小的修正（见习题 7.1）。泡利顺磁性是弱
效应，在多数温度下远小于绝缘体中因居里定律出现的顺磁性。原因：顺磁绝缘体中每个
磁性原子上至少有一个电子作贡献；金属中只有靠近费米面的电子起作用。多数金属顺磁
磁化率之小曾是个谜，直到泡利指出它是电子服从费米--狄拉克统计（而非经典统计）的
后果。

\subsection{向局域行为的过渡}

费米--狄拉克统计的效应，以及泡利顺磁性与局域矩行为之间的过渡，可以通过重新推导
式 7.15 来说明。严格地，每种自旋态的单位体积电子数应写为
\begin{equation*}
\begin{aligned}
n_{\uparrow} &= \frac{1}{2}\int_{0}^{\infty}g(E+\mu_{\mathrm{B}}B)f(E)\,\mathrm{d}E\\
n_{\downarrow} &= \frac{1}{2}\int_{0}^{\infty}g(E-\mu_{\mathrm{B}}B)f(E)\,\mathrm{d}E.
\end{aligned}\tag{7.17}
\end{equation*}
于是小 B 时磁化强度 $M=\mu_{\mathrm{B}}(n_{\uparrow}-n_{\downarrow})$，从而
\begin{align*}
M &\approx \mu_{\mathrm{B}}^{2}B\int_{0}^{\infty}\frac{\mathrm{d}g}{\mathrm{d}E}f(E)\,\mathrm{d}E\\
&= \mu_{\mathrm{B}}^{2}B\left(\left[g(E)f(E)\right]_{0}^{\infty} - \int_{0}^{\infty}\frac{\mathrm{d}f}{\mathrm{d}E}g(E)\,\mathrm{d}E\right)\tag{7.18}
\end{align*}
第二行由第一行分部积分得到。利用 $g(0)=0$ 与 $f(\infty)=0$（见图 7.3），第一项为
零。故
\begin{equation*}
M \approx \mu_{\mathrm{B}}^{2}B\int_{0}^{\infty}\left(-\frac{\mathrm{d}f}{\mathrm{d}E}\right)g(E)\,\mathrm{d}E.\tag{7.19}
\end{equation*}
T=0 的简并极限下，$-\mathrm{d}f/\mathrm{d}E$（阶跃函数的导数）是 $E_{\mathrm{F}}$
处的 delta 函数：
\begin{equation*}
-\frac{\mathrm{d}f}{\mathrm{d}E} = \delta(E-E_{\mathrm{F}}),\tag{7.20}
\end{equation*}
于是回到 $M=\mu_{\mathrm{B}}^{2}Bg(E_{\mathrm{F}})$ 与 $\chi=\mu_0\mu_{\mathrm{B}}^{2}g(E_{\mathrm{F}})$，
与式 7.15 一致。

非简并极限下 $f(E)\approx\mathrm{e}^{-(E-\mu)/k_{\mathrm{B}}T}$，故
\begin{equation*}
-\frac{\mathrm{d}f}{\mathrm{d}E} = \frac{f}{k_{\mathrm{B}}T}\tag{7.21}
\end{equation*}
磁化强度为
\begin{equation*}
M = \frac{\mu_{\mathrm{B}}^{2}B}{k_{\mathrm{B}}T}\int_{0}^{\infty}f(E)g(E)\,\mathrm{d}E = \frac{n\mu_{\mathrm{B}}^{2}B}{k_{\mathrm{B}}T},\tag{7.22}
\end{equation*}
从而
\begin{equation*}
\chi = n\mu_0\mu_{\mathrm{B}}^{2}/k_{\mathrm{B}}T,\tag{7.23}
\end{equation*}
与式 2.28 一致。磁化率等价于单位体积 n 个局域矩（大小 $\mu_{\mathrm{B}}$）。如上
所述，多数金属 $E_{\mathrm{F}}$ 为几 eV，熔点以下一切温度简并极限都成立。然而对
载流子浓度低的材料（如掺杂半导体），$E_{\mathrm{F}}\propto n^{2/3}$ 小得多，可能
达到非简并极限，得到类居里的磁化率。

\subsection{实验技术}

我们导出的金属泡利顺磁性表达式与实验符合尚可，但可以通过修正电子--电子相互作用的
效应加以改进。金属的自旋磁化率可以从 NMR 测量提取：NMR 对传导电子自旋磁矩产生的
场远比对电子轨道运动产生的场（它导致 7.6 节考虑的抗磁效应）敏感。

传导电子自旋与核自旋之间的接触相互作用使核共振频率 $\omega$ 产生小移动
$\Delta\omega$，称为奈特移位（Knight shift）。可以这样理解：设想个别传导电子
在某个核上进进出出；核感受到的净超精细耦合是对所有电子自旋取向的平均。无外场时
电子自旋取向平均为零，净超精细耦合为零；非零静场使电子自旋极化，净超精细耦合
非零。奈特移位 $K=\Delta\omega/\omega$ 因此既正比于核处的传导电子密度（体现对耦合
强度的依赖），又正比于泡利自旋磁化率（体现外场使电子极化的程度）。

超精细相互作用的静态平均造成奈特移位；围绕平均的涨落提供 $T_1$ 弛豫机制（称为
Korringa 弛豫）。主导的 $T_1$ 过程是电子与核（或 $\mu$ 子）自旋的翻转--翻转
（flip--flop）跃迁，其中电子与核塞曼能量之差由传导电子动能的变化补偿。核与传导
电子之间的能量交换非常小，只有费米面附近 $k_{\mathrm{B}}T$ 范围内的电子能够参与——
只有它们附近有空态可供跃迁。因此对简单金属，自旋--晶格弛豫速率 $T_1^{-1}$ 正比于
温度。奈特移位（常以无量纲形式 $\Delta\omega/\omega$ 表示）与 Korringa 弛豫速率
$T_1^{-1}$ 通常由方程
\begin{equation*}
\frac{1}{T_1} \propto T\left(\frac{\Delta\omega}{\omega}\right)^{2},\tag{7.24}
\end{equation*}
联系，称为 Korringa 关系。

\section{自发自旋劈裂能带}

铁中每原子磁矩约 $2.2\mu_{\mathrm{B}}$（见表 5.1）。这个非整数值无法在原子局域矩
的基础上理解——它是能带铁磁性（band ferromagnetism，又称巡游铁磁性）的有力
证据：磁化来源于自发自旋劈裂的能带。本节探讨几个可用于理解某些材料中能带如何
自发自旋劈裂的模型。

\begin{figure}
\centering
\includegraphics[width=0.45\textwidth]{fig7_05.pdf}
\caption{态密度图显示无外加磁场时能带的自发分裂。}
\end{figure}
\bio{Edmund C. Stoner\\（1899--1968）}

在分子场理论中，我们说所有自旋“感受”邻居产生的同一平均交换场 $\lambda M$。
金属中，分子场可以经由泡利顺磁性 $\chi_{\mathrm{P}}$ 磁化电子气；电子气由此产生的
磁化强度 M 又反过来充当分子场。这是“鸡生蛋蛋生鸡”的情形（也称 bootstrap）：
这一正反馈机制能导致自发铁磁性吗？大概能——只要 $\lambda$（给定 M 能得到多少
分子场）与 $\chi_{\mathrm{P}}$（给定分子场能得到多少磁化强度）都足够大。

最好把上述启发式论证做得更严格些！要问的问题是：系统整体通过变为铁磁体能省能量
吗？

先设想：无外场时，把费米面处少量电子从自旋向下带搬到自旋向上带。具体地：取能量从
$E_{\mathrm{F}}-\delta E$ 到 $E_{\mathrm{F}}$ 的自旋向下电子，翻转其自旋放入自旋
向上带，能量从 $E_{\mathrm{F}}$ 到 $E_{\mathrm{F}}+\delta E$（图 7.5）。被搬电子
数为 $g(E_{\mathrm{F}})\delta E/2$，每个能量升高 $\delta E$，总能量变化
$g(E_{\mathrm{F}})\delta E/2\times\delta E$。总动能变化 $\Delta E_{\mathrm{K.E.}}$
为
\begin{equation*}
\Delta E_{\mathrm{K.E.}} = \frac{1}{2}g(E_{\mathrm{F}})(\delta E)^{2}.\tag{7.25}
\end{equation*}
这是能量代价，过程看似不利。然而磁化强度与分子场的相互作用给出可以压倒这一代价
的能量降低。自旋向上的数密度为 $n_{\uparrow}=\frac{1}{2}(n+g(E_{\mathrm{F}})\delta E)$，
向下的为 $n_{\downarrow}=\frac{1}{2}(n-g(E_{\mathrm{F}})\delta E)$。设每个电子磁矩
$1\mu_{\mathrm{B}}$，则磁化强度 $M=\mu_{\mathrm{B}}(n_{\uparrow}-n_{\downarrow})$。
分子场能量为
\begin{equation*}
\Delta E_{\mathrm{P.E.}} = -\int_{0}^{M}\mu_0(\lambda M')\,\mathrm{d}M' = -\frac{1}{2}\mu_0\lambda M^{2} = -\frac{1}{2}\mu_0\mu_{\mathrm{B}}^{2}\lambda(n_{\uparrow}-n_{\downarrow})^{2}.\tag{7.26}
\end{equation*}
写 $U=\mu_0\mu_{\mathrm{B}}^{2}\lambda$（U 量度库仑能\fn{2}{分子场源于交换，交换源于库仑相互作用。更高级的处理中库仑能直接以 $Un_{\uparrow}n_{\downarrow}$ 计入，给出同样结果。}），有
\begin{equation*}
\Delta E_{\mathrm{P.E.}} = -\frac{1}{2}U(g(E_{\mathrm{F}})\delta E)^{2}.\tag{7.27}
\end{equation*}
故总能变化 $\Delta E$ 为
\begin{equation*}
\Delta E = \Delta E_{\mathrm{K.E.}}+\Delta E_{\mathrm{P.E.}} = \frac{1}{2}g(E_{\mathrm{F}})(\delta E)^{2}(1-Ug(E_{\mathrm{F}})).\tag{7.28}
\end{equation*}
$\Delta E<0$ 时自发铁磁性可能，即
\begin{equation*}
Ug(E_{\mathrm{F}}) \geq 1\tag{7.29}
\end{equation*}
称为 Stoner 判据（Stoner criterion）。铁磁不稳定性的这一条件要求库仑效应强、且
费米能处态密度大。若有自发铁磁性，无外场时自旋向上、向下带将劈裂能量 $\Delta$
——即交换劈裂。

若 Stoner 判据不满足，自发铁磁性不会发生，但磁化率可能被改变。把外场与相互作用
的效应都计入即可轻松算出：能量移动 $\delta E$ 产生的磁化强度就是
$M=\mu_{\mathrm{B}}(N_{\uparrow}-N_{\downarrow})=2\mu_{\mathrm{B}}g(E_{\mathrm{F}})\delta E$。于是
\begin{align*}
\Delta E &= \frac{1}{2}g(E_{\mathrm{F}})(\delta E)^{2}(1-Ug(E_{\mathrm{F}})) - MB\tag{7.30}\\
&= \frac{M^{2}}{2\mu_{\mathrm{B}}^{2}g(E_{\mathrm{F}})}(1-Ug(E_{\mathrm{F}})) - MB\tag{7.31}
\end{align*}
它在
\begin{equation*}
\frac{M}{\mu_{\mathrm{B}}^{2}g(E_{\mathrm{F}})}(1-Ug(E_{\mathrm{F}})) - B = 0\tag{7.32}
\end{equation*}
时极小，故磁化率 $\chi$ 为
\begin{equation*}
\chi = \frac{M}{H}\approx\frac{\mu_0M}{B} = \frac{\mu_0\mu_{\mathrm{B}}^{2}g(E_{\mathrm{F}})}{1-Ug(E_{\mathrm{F}})} = \frac{\chi_{\mathrm{P}}}{1-Ug(E_{\mathrm{F}})}.\tag{7.33}
\end{equation*}
这比无库仑相互作用的 $\chi_{\mathrm{P}}$ 大一个因子 $(1-Ug(E_{\mathrm{F}}))^{-1}$，
称为 Stoner 增强（Stoner enhancement）。它负责金属 Pd 与 Pt 中测得的增强泡利
磁化率——二者都可视为处于铁磁性边缘的体系：参数 $Ug(E_{\mathrm{F}})$ 足够大使
磁化率显著增强，但还不够大（不够接近 1）以致自发铁磁。

\section{自旋密度泛函理论}

迄今的讨论用自由电子或近自由电子模型。可以用更高级的方法改进，本节讨论其中之一。
真实体系中不能忽略电子间的库仑相互作用以及交换相互作用对电子运动的影响。所有
粒子的位置与运动相互关联——粒子彼此相互作用、运动中互相施力；相互作用导致粒子
之间出现关联。这类关联理论上很难处理，但一种有用而成功的途径是密度泛函理论
（density functional theory）。

该理论认识到：多电子系统的基态能量可以写成电子密度 $n(\mathbf{r})$ 的泛函
（functional）。\fn{3}{函数（function）是把一个数映射为另一个数的规则。例如函数 $f(x)=x^{2}$ 把数 2 映为数 4。泛函（functional）是把整个函数映射为一个数的规则。例如泛函 $F[f]=\int_{-1}^{1}f(x)\,\mathrm{d}x$ 把函数 $f(x)=x^{2}$ 映为数 2/3。}泛函含三部分贡献：动能；系统中带电粒子静电相互作用的库仑能；
以及称为交换关联能（exchange-correlation energy）的项——它囊括所有多体相互
作用。密度泛函理论无须处理波函数 $\psi(\mathbf{r})$，只须考虑电子密度
$n(\mathbf{r})=|\psi(\mathbf{r})|^{2}$，带来相当大的简化。能量泛函取极小导出一个
可用于求基态能量的方程。

已经证明：极小化泛函确实给出精确的基态能量——尽管用的是电子密度而非波函数。

问题在于能量泛函并非完全已知！动能与库仑能写得出来，交换关联能却是未知的。然而，
可以采用已知的均匀电子气（密度恒定的电子系统）中多电子相互作用的结果作为交换
关联能的取值。这称为局域密度近似（local density approximation，LDA）：在电子
密度为 $n(\mathbf{r})$ 的每个空间点 r，电子对多体相互作用的响应，就好像全空间
（即位置 $\mathbf{r}'\neq\mathbf{r}$ 处）的电子密度都取 $n(\mathbf{r})$。空间不同
位置对能量的贡献全部加起来（对所有体积元积分）；各体积元的贡献不同，依赖于局域
电子密度。LDA 对完美金属（电子密度确实均匀）是精确的，但对电子密度剧烈空间变化
的体系则不那么好。\fn{4}{该理论对金属很有效，对局域体系则不然。原因：库仑能代表某个特定电子对电子密度的响应，而密度来自包括该电子在内的所有电子——电子自相互作用不可避免地、然而错误地被计入。对金属这不算太糟——任一电子只是浩瀚传导电子海洋中的一员；对局域体系则是灾难——特定态可能只被一个电子占据。}

把密度泛函理论推广到包括自旋极化效应，称为自旋密度泛函理论（spin-density
functional theory）。磁性体系中它与局域自旋密度近似（local spin-density
approximation，常称 LSDA）联用：交换关联势不仅依赖局域电子密度，还依赖局域自旋
密度（自旋向上与向下电子密度之差）。这一技术可用于对能带结构做现实计算，获得
真实体系自旋密度的定量信息。不过下一节我们退回到相对安全的简单自由电子模型！

\section{朗道能级}

泡利顺磁性是与电子自旋相联系的效应。电子的轨道贡献呢？用如下论证来评估。设自由
电子哈密顿量为
\begin{equation*}
\frac{\hat{\mathbf{p}}^{2}}{2m_{\mathrm{e}}}\psi = -\frac{\hbar^{2}}{2m_{\mathrm{e}}}\nabla^{2}\psi = E\psi,\tag{7.34}
\end{equation*}
波函数为平面波
\begin{equation*}
\psi = \frac{1}{\sqrt{V}}e^{i\mathbf{k}\cdot\mathbf{r}},\tag{7.35}
\end{equation*}
故能量 E 为
\begin{equation*}
E = \frac{\hbar^{2}k_x^{2}}{2m_{\mathrm{e}}} + \frac{\hbar^{2}k_y^{2}}{2m_{\mathrm{e}}} + \frac{\hbar^{2}k_z^{2}}{2m_{\mathrm{e}}}.\tag{7.36}
\end{equation*}
与前面一样，无磁场时电子态在 k 空间均匀分布、相邻点相距 $2\pi/L$（见图 7.6(a)）。

存在磁场 $\mathbf{B}=(0,0,B)$ 时，动量算符 $\hat{\mathbf{p}}=-\mathrm{i}\hbar\nabla$
须换成 $-\mathrm{i}\hbar\nabla+e\mathbf{A}$，磁矢势的有用选取是
$\mathbf{A}=(0,Bx,0)$。于是式 7.34 变为
\begin{equation*}
\left(\frac{\hat{p}_x^{2}}{2m_{\mathrm{e}}} + \frac{(\hat{p}_y+eB\hat{x})^{2}}{2m_{\mathrm{e}}} + \frac{\hat{p}_z^{2}}{2m_{\mathrm{e}}}\right)\psi = E\psi.\tag{7.37}
\end{equation*}
由于表达式含算符 $\hat{x}$、$\hat{p}_x$、$\hat{p}_y$、$\hat{p}_z$，波函数在 $y$
与 $z$ 方向仍是平面波：$\hat{p}_y\psi=\hbar k_y\psi$、$\hat{p}_z\psi=\hbar k_z\psi$。
于是式 7.37 可整理（用分离变量）为
\begin{equation*}
\left(\left[\frac{\hat{p}_x^{2}}{2m_{\mathrm{e}}} + \frac{1}{2}m_{\mathrm{e}}\omega_{\mathrm{c}}^{2}(x-x_0)^{2}\right] + \frac{\hbar^{2}k_z^{2}}{2m_{\mathrm{e}}}\right)\psi = E\psi,\tag{7.38}
\end{equation*}
$x_0=-\hbar k_y/eB$，且
\begin{equation*}
\omega_{\mathrm{c}} = \frac{eB}{m_{\mathrm{e}}}\tag{7.39}
\end{equation*}
是回旋频率。式 7.38 方括号内是一维谐振子的哈密顿量，故能量本征值对应
\begin{equation*}
E = \left(l+\frac{1}{2}\right)\hbar\omega_{\mathrm{c}} + \frac{\hbar^{2}k_z^{2}}{2m_{\mathrm{e}}}\tag{7.40}
\end{equation*}
l 为整数。能量本征函数是 y、z 方向平面波与 x 方向一维谐振子波函数之积。

在本处理中，$x_0$ 是 x 方向谐振动的“中心”，随 $k_y$ 取不同值。\fn{5}{若选 $\mathbf{A}=(-By,0,0)$，得到的参数 $y_0$ 正比于 $k_x$，是 y 方向谐振动的中心。可以证明对应 $x_0$ 与 $y_0$ 的算符不对易，故 $x_0$ 与 $y_0$ 不能同时取确定值。}处理这一问题的另一个有用方法是选取反映此情形固有的柱对称性的规范
（习题 7.2 这样做），得到与式 7.40 相同的能量本征值，但本征函数是 z 方向平面波与
$\sqrt{k_x^{2}+k_y^{2}}$ 的函数之积。无论用哪种规范求解，允许的态都可以看作位于
平行于 $k_z$ 的朗道管（Landau tubes）上（见图 7.7）。磁场的施加使电子形成朗道
能级（Landau levels，把态按能量看时即如此），每个朗道能级可用量子数 l 标记
（见图 7.6(b)）。朗道能级的简并度可以这样算：k 空间中相邻两个朗道能级之间的
面积（见图 7.6(b)）为
\begin{equation*}
\pi k_{l+1}^{2}-\pi k_l^{2} = \frac{2m_{\mathrm{e}}\pi}{\hbar^{2}}\left(\left(l+1+\frac{1}{2}\right)\hbar\omega_{\mathrm{c}}-\left(l+\frac{1}{2}\right)\hbar\omega_{\mathrm{c}}\right) = \frac{2m_{\mathrm{e}}\pi\omega_{\mathrm{c}}}{\hbar}\tag{7.41}
\end{equation*}
由于每个态被两个电子占据、一个态在 $k_xk_y$ 平面占面积 $(2\pi/L)^{2}$，一个朗道
能级的简并度 p 为
\begin{equation*}
p = \frac{4m_{\mathrm{e}}\pi\omega_{\mathrm{c}}}{\hbar(2\pi/L)^{2}} = \frac{m_{\mathrm{e}}L^{2}\omega_{\mathrm{c}}}{\pi\hbar}\tag{7.42}
\end{equation*}
\mnote{注意：若磁矢势选为 $\mathbf{A}=(-By,0,0)$（它给出同样的
$\mathbf{B}=\nabla\times\mathbf{A}=(0,0,B)$），会得到同样的能量本征值，但本征函数
在此情形将是 x、z 方向平面波与 y 方向一维谐振子波函数之积。“规范的选择”（即
写下 $\mathbf{A}$ 的方式）不影响能量本征值，但影响导出的波函数集合。}

\begin{figure}
\centering
\includegraphics[width=0.4\textwidth]{fig7_06.pdf}
\caption{（a）B=0 时电子态在 k 空间均匀分布，相邻点相距 $2\pi/L$。（b）磁场的施加
使允许能量只剩分立值。这些朗道能级态具有简并度——由波矢 $k_{l+1}$ 与 $k_l$ 能级
之间 k 空间的面积算得。}
\end{figure}

\begin{figure}
\centering
\includegraphics[width=0.45\textwidth]{fig7_07.pdf}
\caption{朗道管。}
\end{figure}

\section{朗道抗磁性}

用上一节的结果，现在可以计算电子气的抗磁响应。想法是：加磁场时电子分布分裂成一
系列朗道能级，系统总能量可能变化；能量随场的变化等价于系统的磁化强度。这一现象
称为朗道抗磁性（Landau diamagnetism）。

计算基于如下观察：不同 $k_z$ 的最高被占据朗道能级不同（见图 7.7——费米面在不同
$k_z$ 处切割不同的朗道管）。对特定 $k_z$，定义与垂直于场的电子运动相联系的能量
$E_{\perp}$：
\begin{equation*}
E_{\perp} = E_{\mathrm{F}} - \frac{\hbar^{2}k_z^{2}}{2m_{\mathrm{e}}}.\tag{7.43}
\end{equation*}
无磁场时，对特定 $k_z$，电子可以填满直到最大值 $E_{\perp}$ 的态；加磁场则允许分立
的朗道能级。把最高被占据朗道能级标记为 $E_l$。B=0 问题中的电子可以分解成若干
“块”，每块的平均能量等于相应的朗道能级能量（顶块例外——部分填充）。

定义参数 x：
\begin{equation*}
x = E_{\perp} - E_l\tag{7.44}
\end{equation*}
（见图 7.8）。最高被占据块的平均能量为
\begin{equation*}
\frac{1}{2}\left[E_{\perp}+E_l-\frac{\hbar\omega_{\mathrm{c}}}{2}\right] = E_{\perp}-\frac{x}{2}-\frac{\hbar\omega_{\mathrm{c}}}{4}\tag{7.45}
\end{equation*}

\begin{figure}
\centering
\includegraphics[width=0.45\textwidth]{fig7_08.pdf}
\caption{（a）B=0 时电子可以填满直到最大值 $E_{\perp}$（由式 7.43 定义）的态；
T=0 被填满的态以阴影示出。文中所述“最高被占据块”标记为 $E_l$；参数
$x=E_{\perp}-E_l$ 量度该块的中点（点划线）与最高被占据态之间的能量。（b）磁场只
允许分立的朗道能级，能量为 $\ldots E_{l-2},E_{l-1},E_l,E_{l+1}\ldots$；T=0 被占据
的朗道能级以粗线示出，最高被占据朗道能级标记为 $E_l$。}
\end{figure}

占据数为
\begin{equation*}
\frac{p}{\hbar\omega_{\mathrm{c}}}\left[E_{\perp}-\left(E_l-\frac{\hbar\omega_{\mathrm{c}}}{2}\right)\right] = \frac{p}{2}\left[1+\frac{2x}{\hbar\omega_{\mathrm{c}}}\right].\tag{7.46}
\end{equation*}
两者相乘给出顶块的能量。这是对特定 $k_z$ 求值的，但必须对费米面上所有电子（即所有
不同 $k_z$ 值）平均。这可以通过对 x 的一切可能值平均实现——设 x 在
$-\hbar\omega_{\mathrm{c}}/2$ 与 $\hbar\omega_{\mathrm{c}}/2$ 之间均匀分布：
\begin{equation*}
\langle x\rangle = \int_{-\hbar\omega_{\mathrm{c}}/2}^{\hbar\omega_{\mathrm{c}}/2}x\,\mathrm{d}x = 0,\tag{7.47}
\end{equation*}
及
\begin{equation*}
\langle x^{2}\rangle = \int_{-\hbar\omega_{\mathrm{c}}/2}^{\hbar\omega_{\mathrm{c}}/2}x^{2}\,\mathrm{d}x = \frac{\hbar^{2}\omega_{\mathrm{c}}^{2}}{12}.\tag{7.48}
\end{equation*}
由此得与垂直于场的运动相联系的顶块（标记 $E_l$ 的那块）平均能量为
\begin{equation*}
\frac{p}{2}\left[E_{\perp}-\frac{\hbar\omega_{\mathrm{c}}}{3}\right]\tag{7.49}
\end{equation*}
总平均能量为（用式 7.43）
\begin{equation*}
\frac{p}{2}\left[E_{\mathrm{F}}-\frac{\hbar\omega_{\mathrm{c}}}{3}\right].\tag{7.50}
\end{equation*}
磁场非零时：若 $0<x<\hbar\omega_{\mathrm{c}}/2$，顶块被 p 个能量 $E_l$ 的电子填满；
若 $-\hbar\omega_{\mathrm{c}}/2<x<0$，顶块全空。（注意此计算针对特定 $k_z$；态在
不同 $k_z$ 之间重新分配。）于是顶块的平均能量为
\begin{equation*}
E = \begin{cases} p(E_{\perp}-x) & \text{若 } x>0\\ 0 & \text{若 } x<0 \end{cases}\tag{7.51}
\end{equation*}
故与垂直于场的运动相联系的平均能量为
\begin{equation*}
\frac{p}{2}\left[E_{\perp}-\frac{\hbar\omega_{\mathrm{c}}}{4}\right]\tag{7.52}
\end{equation*}
总平均能量（再次用式 7.43）为
\begin{equation*}
\frac{p}{2}\left[E_{\mathrm{F}}-\frac{\hbar\omega_{\mathrm{c}}}{4}\right]\tag{7.53}
\end{equation*}
计入 $k_{\mathrm{F}}L/\pi$ 因子（$k_z$ 方向态的简并度），$B\neq0$ 与 B=0 的总能量
之差为
\begin{align*}
\Delta U &= U(B\neq0)-U(B=0)\tag{7.54}\\
&= \frac{pk_{\mathrm{F}}L}{2\pi}\left[E_{\mathrm{F}}-\frac{\hbar\omega_{\mathrm{c}}}{4}\right] - \frac{pk_{\mathrm{F}}L}{2\pi}\left[E_{\mathrm{F}}-\frac{\hbar\omega_{\mathrm{c}}}{3}\right]\tag{7.55}\\
&= \frac{pk_{\mathrm{F}}L\hbar\omega_{\mathrm{c}}}{24\pi}.\tag{7.56}
\end{align*}
用简并度 p 的表达式（式 7.42）：
\begin{equation*}
\Delta U = \frac{Vk_{\mathrm{F}}e^{2}B^{2}}{24\pi^{2}m_{\mathrm{e}}}.\tag{7.57}
\end{equation*}
故磁化率为
\begin{align*}
\chi_{\mathrm{L}} &= -\frac{\mu_0}{V}\frac{\partial^{2}\Delta U}{\partial B^{2}}\tag{7.58}\\
&= -\frac{\mu_0k_{\mathrm{F}}e^{2}}{12\pi^{2}m_{\mathrm{e}}}\tag{7.59}\\
&= -\frac{1}{3}\mu_0\left(\frac{e\hbar}{2m_{\mathrm{e}}}\right)^{2}\left[\frac{m_{\mathrm{e}}k_{\mathrm{F}}}{\pi^{2}\hbar^{2}}\right]\tag{7.60}\\
&= -\frac{1}{3}\mu_0\mu_{\mathrm{B}}^{2}g(E_{\mathrm{F}}).\tag{7.61}
\end{align*}
这就是朗道抗磁性，它与泡利顺磁性 $\chi_{\mathrm{P}}$ 的关系为
\begin{equation*}
\chi_{\mathrm{L}} = -\frac{\chi_{\mathrm{P}}}{3}.\tag{7.62}
\end{equation*}
由此或许会得出结论：一切金属都是顺磁的，因为（正的）泡利顺磁性是（负的）朗道
抗磁性的三倍。这未免轻率——我们用了自由电子模型、忽略了能带结构效应。若电子有
有效质量 $m^{*}$，费米能处态密度 $g(E_{\mathrm{F}})$ 将增强 $m^{*}/m_{\mathrm{e}}$
倍，泡利顺磁性与朗道抗磁性都增强 $m^{*}/m_{\mathrm{e}}$ 倍。但式 7.61 中最终结果
用了 $\mu_{\mathrm{B}}=e\hbar/2m_{\mathrm{e}}$；玻尔磁子是固定常数、不依赖有效
质量。这意味着朗道抗磁性还须再乘 $(m_{\mathrm{e}}/m^{*})^{2}$。净结果：
\begin{equation*}
\chi_{\mathrm{L}} = -\left(\frac{m_{\mathrm{e}}}{m^{*}}\right)^{2}\frac{\chi_{\mathrm{P}}}{3},\tag{7.63}
\end{equation*}
金属的总磁化率为
\begin{equation*}
\chi = \chi_{\mathrm{P}}\left[1-\frac{1}{3}\left(\frac{m_{\mathrm{e}}}{m^{*}}\right)^{2}\right].\tag{7.64}
\end{equation*}
多数金属 $m^{*}\sim m_{\mathrm{e}}$，故为顺磁。若 $m^{*}<m_{\mathrm{e}}/\sqrt{3}$，
净磁化率变负——这解释了铋的强抗磁性（$m^{*}\sim0.01m_{\mathrm{e}}$）。

迄今忽略的另一个效应是离子芯中束缚电子的抗磁性。原子序数增大时该效应变大（贡献
的电子更多），但仍通常远小于我们讨论的传导电子效应。

费米面上的电子运动由方程
\begin{equation*}
\hbar\dot{\mathbf{k}} = -e\mathbf{v}\times\mathbf{B},\tag{7.65}
\end{equation*}
描述，k 是费米面上的电子波矢，$\mathbf{v}=\hbar^{-1}\partial E(\mathbf{K})/\partial\mathbf{k}$
是电子速度，总是垂直于费米面。故电子运动垂直于 B 且保持在费米面上。\fn{6}{这是因为在时间 dt 内 k 的变化 $\mathrm{d}\mathbf{k}\propto(\partial E(\mathbf{K})/\partial\mathbf{k})\times\mathbf{B}$，既垂直于 B 又在费米面平面内。}磁化率的抗磁轨道贡献来自费米面上的闭合轨道。如图 7.9 所示，开放轨道也可能出现。
简单自由电子球中只可能有闭合轨道；真实材料中开、闭轨道皆可能。

\begin{figure}
\centering
\includegraphics[width=0.5\textwidth]{fig7_09.pdf}
\caption{（a）闭合轨道与（b）开放轨道的例子。轨道上的矢量表示电子的实空间速度
$\mathbf{v}=\hbar^{-1}\partial E(\mathbf{K})/\partial\mathbf{k}$——垂直于费米面的
矢量。磁场 B 如（a）取向时得到闭合轨道；B 沿 x 方向时，同一费米面上可以出现开放
轨道，如（b）所示。}
\end{figure}
\bio{P.~M. van Alphen（1906--1967）\\W.~J. de Haas（1878--1960）\\L.~V. Shubnikov（1901--1945）}

磁场增大、朗道能级突破费米面时，性质发生有趣的振荡变化。性质之所以变化，是因为
费米能处的态密度可以作为磁场的函数振荡。当朗道管扫过费米面的极值截面（对费米球
即赤道截面）时效应最大。态密度的振荡导致振荡磁化强度——这就是德哈斯--范阿尔芬
效应（de Haas--van Alphen effect），可用来测量费米面的截面积。电阻上也有类似
效应，称为舒布尼科夫--德哈斯效应（Shubnikov--de Haas effect）。两种效应都给出
性质随 1/B 的振荡。振荡周期按
\begin{equation*}
\Delta\left(\frac{1}{B}\right) = \frac{2\pi e}{\hbar A_{\mathrm{ext}}}.\tag{7.66}
\end{equation*}
给出垂直于磁场的费米面最大与最小截面积 $A_{\mathrm{ext}}$。

\section{电子气的磁性}

泡利顺磁性描述电子气对均匀外场的顺磁响应。若外场随空间变化呢？电子气将按波矢依赖
的磁化率响应这一微扰——现在就来计算（暂且忽略抗磁响应）。

\subsection{电子气的顺磁响应}

空间变化的磁场
\begin{equation*}
\mathbf{H}(\mathbf{r}) = \mathbf{H}_{\mathbf{q}}\cos\mathbf{q}\cdot\mathbf{r}\tag{7.67}
\end{equation*}
提供微扰 $\hat{\mathcal{H}}'$：
\begin{equation*}
\hat{\mathcal{H}}' = \pm\frac{g\mu_0\mu_{\mathrm{B}}}{2}|\mathbf{H}_{\mathbf{q}}|\cos\mathbf{q}\cdot\mathbf{r}\tag{7.68}
\end{equation*}
± 指电子自旋。平面波态
\begin{equation*}
\psi_{\mathbf{k}\pm}(\mathbf{r}) = \frac{1}{\sqrt{V}}e^{i\mathbf{k}\cdot\mathbf{r}}|\pm\rangle\tag{7.69}
\end{equation*}
被微弱地微扰到 $\psi_{(\mathbf{k}+\mathbf{q})\pm}(\mathbf{r})$ 与
$\psi_{(\mathbf{k}-\mathbf{q})\pm}(\mathbf{r})$ 态。这些态的幅度可用一阶微扰论计算
（见附录 C）。受扰波函数为
\begin{equation*}
\psi_{\mathbf{k}\pm}(\mathbf{r}) = \frac{1}{\sqrt{V}}\left(\mathrm{e}^{i\mathbf{k}\cdot\mathbf{r}} \pm \frac{g\mu_0\mu_{\mathrm{B}}H_{\mathbf{q}}}{4}\left[\frac{\mathrm{e}^{\mathrm{i}(\mathbf{k}+\mathbf{q})\cdot\mathbf{r}}}{E_{\mathbf{k}+\mathbf{q}}-E_{\mathbf{k}}} + \frac{\mathrm{e}^{\mathrm{i}(\mathbf{k}-\mathbf{q})\cdot\mathbf{r}}}{E_{\mathbf{k}-\mathbf{q}}-E_{\mathbf{k}}}\right]\right)|\pm\rangle\tag{7.70}
\end{equation*}
$E_k=\hbar^{2}k^{2}/2m_{\mathrm{e}}$。只保留 $H_{\mathbf{q}}$ 的领头阶：
\begin{equation*}
|\psi_{\mathbf{k}\pm}(\mathbf{r})|^{2} = \frac{1}{V}\left(1 \pm \frac{g\mu_0\mu_{\mathrm{B}}H_{\mathbf{q}}m_{\mathrm{e}}}{\hbar^{2}}\times\left[\frac{1}{(\mathbf{k}+\mathbf{q})^{2}-k^{2}} + \frac{1}{(\mathbf{k}-\mathbf{q})^{2}-k^{2}}\right]\cos\mathbf{q}\cdot\mathbf{r}\right).\tag{7.71}
\end{equation*}
受扰波函数导致的磁化强度 $M(\mathbf{r})$ 为
\begin{align*}
M(\mathbf{r}) &= \frac{g\mu_0\mu_{\mathrm{B}}}{2}(|\psi_{k+}(\mathbf{r})|^{2}-|\psi_{k-}(\mathbf{r})|^{2})\\
&= \frac{g^{2}\mu_0\mu_{\mathrm{B}}^{2}m_{\mathrm{e}}H_{\mathbf{q}}\cos\mathbf{q}\cdot\mathbf{r}}{\hbar^{2}V}\left[\frac{1}{(\mathbf{k}+\mathbf{q})^{2}-k^{2}} + \frac{1}{(\mathbf{k}-\mathbf{q})^{2}-k^{2}}\right]\tag{7.72}
\end{align*}
要计算电子气整体的磁响应，须把式 7.72 对费米球内所有电子求和。电子气的空间变化
磁化强度 $M(\mathbf{r})$ 成为
\begin{equation*}
M(\mathbf{r}) = M_{\mathbf{q}}\cos\mathbf{q}\cdot\mathbf{r}\tag{7.73}
\end{equation*}
$M_{\mathbf{q}}$ 为
\begin{align*}
M_{\mathbf{q}} &= \frac{g^{2}\mu_0\mu_{\mathrm{B}}^{2}m_{\mathrm{e}}H_{\mathbf{q}}}{\hbar^{2}V}\int_{|\mathbf{k}|<k_{\mathrm{F}}}g(\mathbf{k})\,\mathrm{d}^{3}\mathbf{k}\left[\frac{1}{(\mathbf{k}+\mathbf{q})^{2}-k^{2}} + \frac{1}{(\mathbf{k}-\mathbf{q})^{2}-k^{2}}\right]\tag{7.74}\\
&= \frac{k_{\mathrm{F}}m_{\mathrm{e}}g^{2}\mu_0\mu_{\mathrm{B}}^{2}H_{\mathbf{q}}}{\pi^{2}\hbar^{2}}\left[1+\frac{4k_{\mathrm{F}}^{2}-q^{2}}{4k_{\mathrm{F}}q}\log\left|\frac{q+2k_{\mathrm{F}}}{q-2k_{\mathrm{F}}}\right|\right]\tag{7.75}
\end{align*}
其中态密度 $g(|k|)=Vk^{2}/\pi^{2}$，并用到积分
\begin{align*}
\int_{|\mathbf{k}|<k_{\mathrm{F}}}\frac{\mathrm{d}^{3}\mathbf{k}}{(\mathbf{k}+\mathbf{q})^{2}-k^{2}} &= \int_{0}^{k_{\mathrm{F}}}2\pi k^{2}\mathrm{d}k\int_{0}^{\pi}\frac{\sin\theta\,\mathrm{d}\theta}{q(q-2k\cos\theta)}\\
&= \int_{0}^{k_{\mathrm{F}}}\frac{\pi k\,\mathrm{d}k}{q}\log\frac{q+2k}{|q-2k|}\\
&= \frac{\pi k_{\mathrm{F}}}{2}\left[1+\frac{4k_{\mathrm{F}}^{2}-q^{2}}{4k_{\mathrm{F}}q}\log\left|\frac{q+2k_{\mathrm{F}}}{q-2k_{\mathrm{F}}}\right|\right]\tag{7.76}
\end{align*}
用方程
\begin{equation*}
\chi_{\mathbf{q}} = M_{\mathbf{q}}/H_{\mathbf{q}}\tag{7.77}
\end{equation*}
（下面将证明其合理）与式 7.10 的 $g(E_{\mathrm{F}})$ 表达式，波矢依赖磁化率
$\chi_{\mathbf{q}}$ 为
\begin{equation*}
\chi_{\mathbf{q}} = \chi_{\mathrm{P}}f\left(\frac{q}{2k_{\mathrm{F}}}\right)\tag{7.78}
\end{equation*}
函数 $f(x)$ 为
\begin{equation*}
f(x) = \frac{1}{2}\left(1+\frac{1-x^{2}}{2x}\log\left|\frac{x+1}{x-1}\right|\right).\tag{7.79}
\end{equation*}
$\chi_{\mathbf{q}}$ 绘于图 7.10。q=0 时 $f(0)=1$，故 $\chi_{q=0}=\chi_{\mathrm{P}}$
——结果令人欣慰地退化为 7.2 节导出的泡利顺磁性。曲线在 $q=2k_{\mathrm{F}}$ 处的
扭折源于费米面的存在（费米球直径为 $2k_{\mathrm{F}}$）。随着维数从三维降低，这一
扭折逐渐加剧——本章后面讨论。

\begin{figure}
\centering
\includegraphics[width=0.5\textwidth]{fig7_10.pdf}
\caption{波矢依赖的顺磁磁化率。q=0 时取泡利顺磁磁化率的值。}
\end{figure}

一般地，波矢依赖磁场 $H_{\mathbf{q}}$ 可由 $\mathbf{H}(\mathbf{r})$ 的傅里叶变换
定义：
\begin{equation*}
\mathbf{H}_{\mathbf{q}} = \int\mathrm{d}^{3}\mathbf{r}\,\mathbf{H}(\mathbf{r})e^{-i\mathbf{q}\cdot\mathbf{r}},\qquad \mathbf{H}(\mathbf{r}) = \frac{1}{(2\pi)^{3}}\int\mathrm{d}^{3}\mathbf{q}\,\mathbf{H}_{\mathbf{q}}e^{i\mathbf{q}\cdot\mathbf{r}}.\tag{7.80}
\end{equation*}
同样，波矢依赖磁化强度 $M_{\mathbf{q}}$ 定义为
\begin{equation*}
\mathbf{M}_{\mathbf{q}} = \int\mathrm{d}^{3}\mathbf{r}\,\mathbf{M}(\mathbf{r})e^{-i\mathbf{q}\cdot\mathbf{r}},\qquad \mathbf{M}(\mathbf{r}) = \frac{1}{(2\pi)^{3}}\int\mathrm{d}^{3}\mathbf{q}\,\mathbf{M}_{\mathbf{q}}e^{i\mathbf{q}\cdot\mathbf{r}},\tag{7.81}
\end{equation*}
波矢依赖磁化率 $\chi_{\mathbf{q}}$ 类似地
\begin{equation*}
\chi_{\mathbf{q}} = \int\mathrm{d}^{3}\mathbf{r}\,\chi(\mathbf{r})\mathrm{e}^{-\mathrm{i}\mathbf{q}\cdot\mathbf{r}},\qquad \chi(\mathbf{r}) = \frac{1}{(2\pi)^{3}}\int\mathrm{d}^{3}\mathbf{q}\,\chi_{\mathbf{q}}\mathrm{e}^{\mathrm{i}\mathbf{q}\cdot\mathbf{r}},\tag{7.82}
\end{equation*}
式 7.80、7.81、7.82 都同时给出了逆变换。波矢依赖磁化率提供了计算电子气磁化强度的
手段。电子气中任意点 r 的磁化强度 $\mathbf{M}(\mathbf{r})$ 与磁场 $\mathbf{H}(\mathbf{r})$
的关系为
\begin{equation*}
\mathbf{M}(\mathbf{r}) = \int\mathrm{d}^{3}\mathbf{r}'\,\chi(\mathbf{r}-\mathbf{r}')\mathbf{H}(\mathbf{r}')\tag{7.83}
\end{equation*}
这是一个卷积。故 $\mathbf{M}(\mathbf{r})$ 不仅响应 r 处的场，还响应附近值的加权
平均。卷积定理意味着相应波矢依赖量之间的关系为
\begin{equation*}
\mathbf{M}_{\mathbf{q}} = \chi_{\mathbf{q}}\mathbf{H}_{\mathbf{q}},\tag{7.84}
\end{equation*}
与我们上面用的式 7.77 相同。

\subsection{电子气的抗磁响应}

波矢依赖的抗磁磁化率也可以计算，只是相当复杂。这里只引用结果：
\begin{equation*}
\chi_{\mathbf{q}} = \chi_{\mathrm{L}}\frac{3k_{\mathrm{F}}^{2}}{2q^{2}}\left[1+\frac{4k_{\mathrm{F}}^{2}}{q^{2}} - \frac{k_{\mathrm{F}}}{q}\left(1-\frac{q^{2}}{4k_{\mathrm{F}}^{2}}\right)^{2}\log\left|\frac{q+2k_{\mathrm{F}}}{q-2k_{\mathrm{F}}}\right|\right].\tag{7.85}
\end{equation*}
抗磁磁化率 $\chi_{\mathbf{q}}$ 绘于图 7.11。q=0 时表达式退化为
$\chi_{q=0}=\chi_{\mathrm{L}}=-\chi_{\mathrm{P}}/3$。式 7.85 在 $q=2k_{\mathrm{F}}$
处同样有奇异性——费米面存在之故。

\begin{figure}
\centering
\includegraphics[width=0.5\textwidth]{fig7_11.pdf}
\caption{波矢依赖的抗磁磁化率。q=0 时取朗道抗磁磁化率的值。}
\end{figure}

\subsection{RKKY 相互作用}

波矢依赖的顺磁磁化率对 delta 函数微扰系统的实空间行为有重要后果。delta 函数微扰
$\mathbf{H}(\mathbf{r})=\delta(\mathbf{r})\mathbf{H}$ 可以分解成所有可能空间频率
之和：
\begin{equation*}
\mathbf{H}(\mathbf{r}) = \delta(\mathbf{r})\mathbf{H} = \frac{1}{(2\pi)^{3}}\int\mathbf{H}_{\mathbf{q}}e^{i\mathbf{q}\cdot\mathbf{r}}\,\mathrm{d}^{3}q\tag{7.86}
\end{equation*}
$H_{\mathbf{q}}=H$。delta 函数因此可以看作一切频率之和——包括波长极长与极短的
振荡。若 $\chi_{\mathbf{q}}=\chi$ 与 q 无关，则 $M_{\mathbf{q}}=\chi H_{\mathbf{q}}$，
且
\begin{equation*}
\mathbf{M} = \frac{1}{(2\pi)^{3}}\int\mathbf{M}_{\mathbf{q}}e^{i\mathbf{q}\cdot\mathbf{r}}\,\mathrm{d}^{3}q = \chi\delta(\mathbf{r})\mathbf{H},\tag{7.87}
\end{equation*}
即磁化强度也是 delta 函数：电子气完全响应微扰、完美地跟随其行为。（此时
$\chi(\mathbf{r}-\mathbf{r}')=\chi\delta(\mathbf{r}-\mathbf{r}')$。）

然而电子气的磁化率依赖 q。由于费米面的存在，波矢高于 $q=2k_{\mathrm{F}}$ 时电子气
的响应能力大大降低。因此它虽然能跟上 delta 函数微扰的长波长分量，短波长成分
（波长短于约 $\pi/k_{\mathrm{F}}$）被衰减——电子气难以响应如此高的空间频率。
这在磁化强度中产生空间“振铃”效应，类比单缝衍射图样中的振荡。由此产生的磁化
强度具有空间振荡。

可以这样计算。实空间磁化率为
\begin{align*}
\chi(\mathbf{r}) &= \frac{1}{(2\pi)^{3}}\int\mathrm{d}^{3}\mathbf{q}\,\chi_{\mathbf{q}}\mathrm{e}^{\mathrm{i}\mathbf{q}\cdot\mathbf{r}}\\
&= \frac{1}{(2\pi)^{3}}\int\mathrm{d}^{3}\mathbf{q}\,\frac{\chi_{\mathrm{P}}}{2}\left(1+\frac{4k_{\mathrm{F}}^{2}-q^{2}}{4k_{\mathrm{F}}q}\log\left|\frac{q+2k_{\mathrm{F}}}{q-2k_{\mathrm{F}}}\right|\right)\mathrm{e}^{\mathrm{i}\mathbf{q}\cdot\mathbf{r}}\\
&= \frac{2k_{\mathrm{F}}^{3}\chi_{\mathrm{P}}}{\pi}F(2k_{\mathrm{F}}r)\tag{7.88}
\end{align*}
$r=|\mathbf{r}|$，函数 $F(x)$ 为
\begin{equation*}
F(x) = \frac{-x\cos x+\sin x}{x^{4}}\tag{7.89}
\end{equation*}
绘于图 7.12。这就是 RKKY 相互作用（4.2.4 节已介绍）。式 7.88 的积分绝非平凡，
习题 7.5 将更详细地考察。磁化率（因而 delta 函数微扰所致的磁化强度）在大距离
（$r\gg k_{\mathrm{F}}^{-1}$）正比于 $\cos(2k_{\mathrm{F}}r)/r^{3}$ 且振荡；小 r
处磁化率发散——这是我们假定的 delta 函数微扰的后果。RKKY 相互作用最初是为理解
局域核磁矩对电子气的效应而提出的：核磁矩几乎是最接近 delta 函数的微扰，但它并非
无限局域，而是摊在（尽管极小的）核体积上。RKKY 相互作用为金属中局域电子磁矩之间
的磁耦合提供重要机制：一个磁矩产生电子气的振荡磁化强度，后者可与第二个磁矩相互
作用。依两磁矩间距不同，相互作用可铁磁可反铁磁。

非磁情形也观察到类似现象——电子气对 delta 函数电荷的响应。这与金属合金中局域
电荷的研究有关：电荷密度中的振荡由类似机制产生，称为 Friedel 振荡。

\begin{figure}
\centering
\includegraphics[width=0.5\textwidth]{fig7_12.pdf}
\caption{式 7.89 的函数 $F(x)$——描述 delta 函数场产生的实空间磁化强度。}
\end{figure}

\section{电子气中的激发}

电子气还包含丰富的激发谱。优雅的处理是计算 $\chi(\mathbf{q},\omega)$——依赖
q 与 $\omega$ 的磁化率——但本书不采用此法。

在无磁场时能带就已自旋劈裂的巡游铁磁体中，有两类激发。其一是熟悉的自旋波激发，
遵循常规色散关系：我们 6.6 节的自旋波处理基于海森伯模型、适用于局域矩系统，但
事实证明金属体系中也能导出自旋波。其二是电子--空穴激发的连续谱：电子从一个自旋
劈裂带的已填满态转移到另一自旋劈裂带的空态。后者称为 Stoner 激发（Stoner
excitations）。Stoner 激发中，波矢 $\mathbf{k}+\mathbf{q}$、自旋向下的电子被
激发到波矢 $\mathbf{k}$、自旋向上的态。激发能量为
\begin{equation*}
\hbar\omega = E_{\mathbf{k}+\mathbf{q}} - E_{\mathbf{k}} + \Delta.\tag{7.90}
\end{equation*}
$E_k=\hbar^{2}k^{2}/2m_{\mathrm{e}}$，$\Delta$ 是交换劈裂——翻转一个自旋的能量
代价。结果示于图 7.13。

\begin{figure}
\centering
\includegraphics[width=0.31\textwidth]{fig7_13.pdf}
\caption{电子气中的激发。（a）自旋劈裂带间隔为交换劈裂 $\Delta$。（b）自旋向上与
自旋向下电子的费米面。（c）若存在波矢 $\mathbf{k}$ 的已填满 $\uparrow$ 态与波矢
$\mathbf{k}+\mathbf{q}$ 的空 $\downarrow$ 态，则可以产生 Stoner 激发；阴影区给出
不同 $\mathbf{q}$ 值下 $\mathbf{k}$ 的可能选取。（d）色散关系（$\omega$ 对 $q$
的图）显示自旋波支与 Stoner 激发连续谱。}
\end{figure}

顺磁金属中两带在无磁场时并不自旋劈裂。自旋波激发寿命很短、被强烈阻尼（若视作
阻尼谐振子则处于过阻尼极限）；它们称为顺磁振子（paramagnons），与常规自旋波
激发一样可以用非弹性中子散射研究。

本书迄今用的是两幅截然不同的图景：绝缘材料——磁性联系于原子上局域矩；金属——
磁性联系于完全离域的矩。许多真实材料介于两者之间：磁性联系于介于局域矩与能带铁磁
之间的自旋密度涨落。局域矩涨落在实空间局域，而弱铁磁金属中的自旋涨落可视为在
q 空间局域。

\section{自旋密度波}

Stoner 判据表明：$Ug(E_{\mathrm{F}})$ 足够大时电子气对自发铁磁性不稳定（见
7.3 节），磁化率可被增强 $(1-Ug(E_{\mathrm{F}}))^{-1}$ 倍。事实证明，波矢依赖
磁化率（无库仑相互作用时记作 $\chi_{\mathbf{q}}^{0}=\chi_{\mathrm{P}}f(q/2k_{\mathrm{F}})$）
也被库仑相互作用增强为
\begin{equation*}
\chi_{\mathbf{q}} = \frac{\chi_{\mathrm{P}}f(q/2k_{\mathrm{F}})}{1-Ug(E_{\mathrm{F}})f(q/2k_{\mathrm{F}})} = \frac{\chi_{\mathbf{q}}^{0}}{1-\alpha\chi_{\mathbf{q}}^{0}}\tag{7.91}
\end{equation*}
$\alpha=U/\mu_0\mu_{\mathrm{B}}^{2}$。可能出现：$\chi_{\mathbf{q}}^{0}$ 在某个
q≠0 处取极大。此时若 $\alpha\chi_{\mathbf{q}}^{0}$ 达到 1，参数化为 $\alpha$ 的
相互作用可使磁化率在该 q 值发散——样品中可能自发发展出振荡的静态磁化强度。
q=0 对应铁磁有序；$|q|=\pi/a$ 对应反铁磁有序。一般地，预期出现波矢 q 的螺旋结构
或自旋密度波结构。

\begin{figure}
\centering
\includegraphics[width=0.5\textwidth]{fig7_14.pdf}
\caption{一、二、三维的波矢依赖顺磁磁化率。一维情形响应函数在 q=2k$_{\mathrm{F}}$ 处
发散。所有维数下 q=0 处都取泡利顺磁磁化率的值。}
\end{figure}

我们已对三维电子气求出了波矢依赖磁化率，但对二维与一维可以重复计算（见习题
7.6）。结果示于图 7.14：降低维数时，三维曲线在 $q=2k_{\mathrm{F}}$ 处的“膝”在
二维变成扭折、在一维变成奇异性。一维的峰称为 Kohn 异常（Kohn anomaly），表明
一维电子气对形成波矢 $2k_{\mathrm{F}}$ 的自旋密度波不稳定。这一不稳定与如下事实
相关：波长 $2\pi/q$、波矢 q 的磁化强度周期性调制会在波矢 $\pm q/2$ 处打开能隙
（任何周期势都能做到）。\bio{Walter Kohn\\（1923--　）}若 $q/2=k_{\mathrm{F}}$，能隙恰在费米面处，可以降低总能量——降低电子能量的机会本身就能驱动自旋密度波的形成。

\begin{figure}
\centering
\includegraphics[width=0.45\textwidth]{fig7_15.pdf}
\caption{（a）能带中电子填充直到费米能 $E_{\mathrm{F}}$。（b）波矢 $q=2k_{\mathrm{F}}$
的周期势通过在费米面产生能隙降低总电子能量。费米面打开能隙（从（a）到（b））伴随
金属--绝缘体转变；相应的实空间畸变称为 Peierls 畸变。一维电子气对这种 Peierls
不稳定性是不稳定的。}
\end{figure}
\bio{Rudolf E. Peierls\\（1907--1995）}

波矢 $2k_{\mathrm{F}}$ 的电荷密度波的形成有类似的不稳定性。两种效应都与著名的
Peierls 不稳定性相关：一维原子链的二聚化可以自发发生，因为在费米面附近打开能隙
所节省的电子能量可以超过二聚化的弹性能量代价（见图 7.15）。

三维中自旋密度波无法在费米面所有点产生能隙。有时费米面是这样的：把费米面的一部分
平移矢量 q 后恰好落在另一部分之上——这一现象称为嵌套（nesting）。波矢 q 的自旋
密度波可以在费米面上可嵌套的区域产生能隙。真实金属中经常出现：两块费米面在 k 空间
中被固定波矢 q 近似平移相接。这可以产生磁化率峰与由此的不稳定性：密度波的形成
被称为嵌套了费米面。

若嵌套波矢 q 恰为 $\pi/a$（a 为原子间距），自旋密度波与晶格可公度，产生反铁磁
有序（图 7.16(a)）。但更常见的情形是嵌套波矢 q（等于 $2k_{\mathrm{F}}$）不是
$\pi/a$ 的简单倍数，自旋密度波不可公度（图 7.16(b)）。

\begin{figure}
\centering
\includegraphics[width=0.62\textwidth]{fig7_16.pdf}
\caption{（a）可公度的自旋密度波，波矢 $q=\pi/a$。（b）不可公度自旋密度波。}
\end{figure}

铬具有体心立方结构（立方边长 a），在 $T_{\mathrm{N}}=310$ K 以下发展出自旋密度
波结构：波矢 $\mathbf{q}\approx0.96(2\pi/a)$（故材料“几乎”是反铁磁体），随温度
缓慢变化，方向沿 [100]。极化在 $T<115$ K 时为纵向、$T>115$ K 时为横向。Cr 的
电阻率在略低于 $T_{\mathrm{N}}$ 处上升——自旋密度波在部分费米面上引入能隙、减少
载流子数目所致。若干有机金属中也发现自旋密度波。

\section{近藤效应}

非磁宿主中磁性离子的极稀合金中，若忽略 RKKY 耦合，磁性离子的磁矩可视为独立。
剩下的唯一相互作用是磁矩与传导电子自旋之间。高温下磁矩表现得像自由顺磁矩；低于
特征温度——近藤温度（Kondo temperature）$T_{\mathrm{K}}$——磁矩与传导电子的
相互作用使杂质自旋变得非磁。实际发生的是：传导电子开始在杂质自旋周围形成相反
自旋极化的云，形成准束缚态。传导电子对磁性杂质的这种磁屏蔽过程称为近藤效应
（Kondo effect），有两个深刻的实验后果。其一，磁化强度降到自由矩值以下，磁化率
因此低于居里定律的预期。其二，由于杂质矩与传导电子强烈相互作用，矩的电子散射
截面强烈增强，电阻率出现正比于 $\mathsf{J}\ln T$ 的新项（$\mathsf{J}<0$ 是局域矩
与传导电子之间的交换耦合）。

多数材料的电阻率随温度降低而减小：材料冷却时声子数目减少，而电阻率强烈依赖声子
数目。故通常高温下由正比于 T 的项主导，较低温度下由 $T^{5}$ 主导；最低温度时由
杂质浓度（不依赖温度）决定的常数项变得重要。

近藤效应的新 $\mathsf{J}\ln T\propto-\ln T$ 项（因 $\mathsf{J}<0$）在低温介入，
导致电阻率极小的出现（见图 7.17）：低于极小值温度，电阻率随降温而上升。

\begin{figure}
\centering
\includegraphics[width=0.42\textwidth]{fig7_17.pdf}
\caption{电阻率的温度依赖显示极小。电阻率是正比于 $T^{5}$ 的项与 $\mathsf{J}\ln T\propto-\ln T$ 项之和。}
\end{figure}

\section{哈伯德模型}

金属行为之所以出现，是因为电子通过在整个晶体上离域化而节省动能。然而在某些体系
中，在位库仑能（把两个电子放到同一格点的代价）强到电子无法在晶体中自由移动；
双占据的高昂代价意味着金属行为失效。电子不再能当作自由粒子处理，而电子关联——
库仑效应大时电子必须相互回避的那些特征——变得非常重要。金属行为（由紧束缚模型
中的转移积分或电子带宽的大小参数化）与库仑能（由 U 参数化，有时称 Hubbard-U）
的竞争由哈伯德模型（Hubbard model）表达。

材料中每个格点可以双占据——一个自旋向上电子与一个自旋向下电子。泡利不相容原理
禁止同自旋电子双占据同一格点。半满极限下（电子恰好每格点一个），若每个格点上
坐一个电子且不移动，在位库仑能可以极小。$U=\infty$ 时这就是基态。U 大而非无穷
时，每个格点仍单占据，但若邻居自旋反平行，电子可以短暂拜访邻居的格点——省一点
动能，因此邻居格点上电子反平行有能量好处，系统将是反铁磁体。但电子无法到次近邻
格点——那里自旋相同，泡利原理禁止同自旋双占据。因此系统还是绝缘体。这种 U 大、
每格点一个电子的反铁磁绝缘体称为莫特绝缘体（Mott insulator）。U 小于电子带宽 W
时实现金属行为。交叉发生在 $U\sim W$——存在金属--绝缘体转变。
\bio{Nevill F. Mott\\（1905--1996）}

\section{中子星}

本章以磁性凝聚态物质的一个极端例子结束：中子星。中子星形成于超新星爆发中恒星的
快速坍缩，由高密核物质组成。恒星坍缩中角动量守恒意味着前身星原有的转速大为增加
（星坍缩时转动惯量骤减、转速飙升）。中子星半径典型 $\sim10$ km，转速可达
$10^{3}\,\mathrm{s}^{-1}$。坍缩还伴随磁通压缩，因此中子星的磁场巨大，可达
$\sim10^{8}$ T。作为对比，太阳表面磁场 $\sim10^{-4}$ T，地球表面 $\sim5\times10^{-5}$
T。磁场在极区附近最强，但中子星的磁轴与地球一样相对自转轴倾斜，故随中子星旋转而
进动。变化的磁场经法拉第定律产生大电场，沿磁轴加速电子；电子沿弯曲轨迹被加速而
发出电磁辐射——形成随中子星自转而旋转的紧致辐射束。地球恰好处在束流路径上时，
我们便收到规则的辐射脉冲。因此中子星最初被称为脉冲星（pulsar）——脉动的星。

用“自由中子”模型并作非相对论处理，可以对中子星内部中子的行为做些（很粗糙的）
估计。中子与电子一样是费米子、服从泡利不相容原理。作为中子星的初步模型，可以
按 7.1 节的方法考虑中子填满中子星大小的盒子里的 k 态。中子星内部主要是中子，密度
比金属中传导电子高约 12 个量级；中子质量当然比电子质量高约 3 个量级。由于费米能
正比于 $n^{2/3}/m$，中子星中的费米能应比金属中高约 5 个量级。于是中子星中中子的
费米温度\fn{7}{作为对比，金属中的电子 $E_{\mathrm{F}}\sim1$ eV，故 $T_{\mathrm{F}}\sim10^{4}$ K。}应约 $10^{9}$ K，远高于中子星温度（可达 $\sim10^{6}$ K）——表明中子星中的中子处于简并极限。中子星表面据认为是固体，内部物质可能处于超流态。这类似于
超流氦中氦原子的配对；中子星中凝聚的可能是中子对、也有质子对。中子与质子的配对能
为 1--2 MeV，相当于 $\sim10^{10}$ K——远高于中子星温度，故可预期配对发生。

\section*{延伸阅读}

\begin{itemize}[itemsep=0.2em]
\item 朗道抗磁性的讨论见 A.~Dupr\'{e}, American Journal of Physics, \textbf{49},
34 (1981)；本章的处理改编自该文。
\item 金属能带理论与德哈斯--范阿尔芬效应的清晰描述见 J.~Singleton,《Electronic
properties of solids》, OUP 2001。
\item 朗道能级物理的有用处理见 J.~H.~Davies,《The physics of low-dimensional
semiconductors》第 6 章，CUP 1998。
\item 能带电子磁性的高级论述见 K.~Yosida,《Theory of magnetism》, Springer 1996。
\item 自旋密度泛函理论见 J.~K\"{u}bler,《Theory of itinerant electron magnetism》,
OUP 2000。
\item 金属磁体中涨落的有用信息见 T.~Moriya,《Spin fluctuations in itinerant
electron magnetism》, Springer 1985。
\item 近藤效应见 A.~C.~Hewson,《The Kondo problem to heavy fermions》, CUP 1993。
\item 中子星及相关天体的物理见 S.~L.~Shapiro 与 S.~A.~Teukolsky,《Black holes,
white dwarfs and neutron stars》, Wiley 1983。
\end{itemize}

\section*{习题}

\exer{7.1} 证明泡利顺磁性给出的磁化率可写成类居里形式
\begin{equation*}
\chi_{\mathrm{P}} = \frac{n\mu_0\mu_{\mathrm{B}}^{2}}{k_{\mathrm{B}}T_0}\tag{7.92}
\end{equation*}
常数 $T_0$ 为
\begin{equation*}
T_0 = \frac{2T_{\mathrm{F}}}{3}\tag{7.93}
\end{equation*}
$T_{\mathrm{F}}=E_{\mathrm{F}}/k_{\mathrm{B}}$ 是费米温度。再证明 $T\ll T_{\mathrm{F}}$
时 $\chi_{\mathrm{P}}$ 的温度依赖修正为
\begin{equation*}
\chi_{\mathrm{P}}(T) = \frac{\mu_0\mu_{\mathrm{B}}^{2}}{3k_{\mathrm{B}}T_0}\left(1-\frac{\pi^{2}T^{2}}{12T_{\mathrm{F}}^{2}}\right)\tag{7.94}
\end{equation*}
并估算典型金属在室温下该修正的大小。

\exer{7.2} （a）证明含 N 个电子的非简并电子气的顺磁磁化率，与 N 个轨道运动被淬灭
的独立局域电子的磁化率相同。

（b）考虑导带含 $3\times10^{22}$ 个电子每立方米、有效质量 $m^{*}=0.1m_{\mathrm{e}}$
的半导体。估算磁化率在何温度以下变得不依赖温度。低于该温度时，计算泡利顺磁性与朗道
抗磁性的大小。

\exer{7.3} 存在磁场 $\mathbf{B}=(0,0,B)$ 时，动量算符 $\hat{\mathbf{p}}=-\mathrm{i}\hbar\nabla$
须换成 $-\mathrm{i}\hbar\nabla+e\mathbf{A}$，导致式 7.40 的能量本征值。7.5 节选了
$\mathbf{A}=(0,Bx,0)$。证明改用 $\mathbf{A}=(-By,0,0)$ 可得同样结果。问题也可以
在柱坐标中求解：$A_{\phi}=\frac{1}{2}Br$、$A_z=A_r=0$。证明这一选择对 $\mathbf{B}$
与能量本征值给出相同结果。

\exer{7.4} 证明费米球的傅里叶变换与式 7.89 的 RKKY 函数相关，即
\begin{align*}
\int_{|\mathbf{k}|<k_{\mathrm{F}}}\mathrm{d}^{3}\mathbf{k}\,\mathrm{e}^{\mathrm{i}\mathbf{k}\cdot\mathbf{r}} &= \int_{0}^{k_{\mathrm{F}}}2\pi k^{2}\mathrm{d}k\int_{0}^{\pi}\mathrm{e}^{\mathrm{i}kr\cos\theta}\sin\theta\,\mathrm{d}\theta\\
&= \frac{4\pi}{r^{3}}\left[\sin k_{\mathrm{F}}r - k_{\mathrm{F}}r\cos k_{\mathrm{F}}r\right]\tag{7.95}
\end{align*}

\exer{7.5} 式 7.88 的积分求值有点复杂，本题设计为一步一步来。先证明
\begin{equation*}
\int\mathrm{d}^{3}\mathbf{q}\,\mathrm{e}^{\mathrm{i}\mathbf{q}\cdot\mathbf{r}}\,\varphi(\mathbf{q}) = \frac{2\pi}{\mathrm{i}r}\int_{-\infty}^{\infty}\mathrm{d}q\,q\,\mathrm{e}^{\mathrm{i}qr}\varphi(q),\tag{7.96}
\end{equation*}
$\varphi(q)=\varphi(-q)$ 是变量 q 的函数。进而证明
\begin{align*}
\frac{1}{(2\pi)^{3}}\int\mathrm{d}^{3}\mathbf{q}\, & \frac{\chi_{\mathrm{P}}}{2}\left(1+\frac{4k_{\mathrm{F}}^{2}-q^{2}}{4k_{\mathrm{F}}q}\log\left|\frac{q+2k_{\mathrm{F}}}{q-2k_{\mathrm{F}}}\right|\right)\mathrm{e}^{\mathrm{i}\mathbf{q}\cdot\mathbf{r}}\tag{7.97}\\
&= \frac{\chi_{\mathrm{P}}k_{\mathrm{F}}^{2}}{\mathrm{i}\pi^{2}r}\int_{-\infty}^{\infty}\mathrm{d}x\,x\,\mathrm{e}^{\mathrm{i}yx}f(x),\tag{7.98}
\end{align*}
$x=q/2k_{\mathrm{F}}$、$y=2k_{\mathrm{F}}r$，$f(x)$ 是式 7.79 定义的函数。求此
积分需要借助围道积分的技巧。

考虑两个积分
\begin{align*}
I_1 &= \frac{1}{\mathrm{i}r}\int_{-\infty}^{\infty}\mathrm{d}x\,\frac{x\,\mathrm{e}^{\mathrm{i}yx}}{2}\left(1+\frac{1-x^{2}}{2x}\log\left|\frac{x+1}{x-1}\right|\right)\tag{7.99}\\
I_2 &= \frac{1}{\mathrm{i}r}\int_{-\infty}^{\infty}\mathrm{d}x\,\frac{x\,\mathrm{e}^{\mathrm{i}yx}}{2}\left(1+\frac{1-x^{2}}{2x}\log\frac{x+1}{x-1}\right).\tag{7.100}
\end{align*}
两个被积函数在 $|x|>1$ 时相等。把围道用 x 平面上半平面的无穷大半圆补全并用
Jordan 引理，证明 $I_2$ 为零。于是利用这一结果与一条从 x=−1 到 x=1、在实轴上方
无穷小位移的围道计算 $I_1$，证明
\begin{align*}
I_1 &= \frac{1}{\mathrm{i}r}\int_{-1}^{1}\mathrm{d}x\,\frac{\mathrm{e}^{\mathrm{i}yx}(1-x^{2})}{4}\left(\log\left|\frac{x+1}{x-1}\right|-\log\frac{x+1}{x-1}\right)\\
&= \frac{\pi}{r}\int_{-1}^{1}\mathrm{d}x\,\frac{\mathrm{e}^{\mathrm{i}yx}(1-x^{2})}{4}\\
&= \frac{\pi}{r}\left[\frac{\sin y-y\cos y}{y^{3}}\right]\\
&= 2\pi k_{\mathrm{F}}F(2k_{\mathrm{F}}r),\tag{7.101}
\end{align*}
从而验证式 7.88。

\exer{7.6} （a）证明一、二、三维费米能级处的态密度分别为
\begin{equation*}
\begin{aligned}
g(E_{\mathrm{F}}) &= \frac{n}{2E_{\mathrm{F}}} = \frac{2m}{\hbar^{2}\pi k_{\mathrm{F}}}\\
g(E_{\mathrm{F}}) &= \frac{n}{E_{\mathrm{F}}} = \frac{m}{\hbar^{2}\pi}\\
g(E_{\mathrm{F}}) &= \frac{3n}{2E_{\mathrm{F}}} = \frac{mk_{\mathrm{F}}}{\hbar^{2}\pi^{2}}
\end{aligned}\tag{7.102}
\end{equation*}
（一维、二维、三维）

（b）证明一维电子气的波矢依赖磁化率为
\begin{equation*}
\chi_{\mathbf{q}} = \chi_{\mathrm{P}}\frac{k_{\mathrm{F}}}{q}\log\left|\frac{q+2k_{\mathrm{F}}}{q-2k_{\mathrm{F}}}\right|\tag{7.103}
\end{equation*}
二维为
\begin{equation*}
\chi_{\mathbf{q}} = \begin{cases} \chi_{\mathrm{P}} & q<2k_{\mathrm{F}}\\ \chi_{\mathrm{P}}\left(1-\sqrt{1-(4k_{\mathrm{F}}^{2}/q^{2})}\right) & q>2k_{\mathrm{F}}. \end{cases}\tag{7.104}
\end{equation*}
这些结果绘于图 7.14。

\exer{7.7} 证明 7.8 节描述、图 7.13 绘出的 Stoner 激发连续谱以两条抛物线
\begin{equation*}
\hbar\omega = \Delta + \frac{\hbar^{2}q}{2m_{\mathrm{e}}}(q\pm2k_{\mathrm{F}\uparrow})\tag{7.105}
\end{equation*}
为界，且 $\omega=0$ 的激发包含在
\begin{equation*}
k_{\mathrm{F}\uparrow}-k_{\mathrm{F}\downarrow} \leq q \leq k_{\mathrm{F}\uparrow}+k_{\mathrm{F}\downarrow},\tag{7.106}
\end{equation*}
的 q 范围内，如图 7.13 所示。（提示：利用 $\hbar^{2}k_{\mathrm{F}\uparrow}^{2}/2m_{\mathrm{e}}=\hbar^{2}k_{\mathrm{F}\downarrow}^{2}/2m_{\mathrm{e}}+\Delta$。）

\exer{7.8} 考虑图 7.5 所示的能带劈裂，但现在 $\delta E$ 不是无穷小量。费米能变得
依赖自旋。若电子系统每单位体积有 n 个电子，证明每单位体积自旋向上电子数 $n_{+}$
与自旋向下电子数 $n_{-}$ 为
\begin{align*}
n_{\pm} &= \int_{0}^{E_{\mathrm{F}\pm}}g_{\pm}(E)\,\mathrm{d}E\tag{7.107}\\
&= \int_{0}^{E_{\mathrm{F}\pm}}\frac{3}{4}\frac{n}{E_{\mathrm{F}}^{3/2}}E^{1/2}\,\mathrm{d}E\tag{7.108}\\
&= \frac{n}{2}\left(\frac{E_{\mathrm{F}\pm}}{E_{\mathrm{F}}}\right)^{3/2},\tag{7.109}
\end{align*}
$E_{\mathrm{F}}$ 是无交换劈裂时的费米能。

用 $n_{+}+n_{-}=n$ 并定义 $x=(n_{+}-n_{-})/n$，证明
\begin{equation*}
\frac{E_{\mathrm{F}\pm}}{E_{\mathrm{F}}} = (1\pm x)^{2/3}.\tag{7.110}
\end{equation*}
动能为
\begin{equation*}
T = \sum_{\sigma=\pm}\int_{0}^{E_{\mathrm{F}\sigma}}E\,g(E)\,\mathrm{d}E.\tag{7.111}
\end{equation*}
证明
\begin{equation*}
T = \frac{3}{10}nE_{\mathrm{F}}\left[(1+x)^{5/3}+(1-x)^{5/3}\right].\tag{7.112}
\end{equation*}
相互作用能可写为
\begin{equation*}
V = -\int_{0}^{M}\mu_0(\lambda M')\,\mathrm{d}M' = -\frac{1}{2}\mu_0\lambda M^{2} = -\frac{1}{2}Un^{2}x^{2}\tag{7.113}
\end{equation*}
其中 $\lambda$、$U=\mu_0\mu_{\mathrm{B}}^{2}\lambda$ 与
$M=\mu_{\mathrm{B}}(n_{+}-n_{-})$ 分别是分子场参数、库仑能与磁化强度。总能量为
\begin{equation*}
E = T+V.\tag{7.114}
\end{equation*}
证明条件 $\mathrm{d}E/\mathrm{d}x=0$ 在
\begin{equation*}
\frac{(1+x)^{2/3}+(1-x)^{2/3}}{x} = \frac{4Ug(E_{\mathrm{F}})}{3}.\tag{7.115}
\end{equation*}
时给出稳定解。（提示：用 $g(E_{\mathrm{F}})=3n/2E_{\mathrm{F}}$。）作图解释这一
结果。

证明当
\begin{equation*}
Ug(E_{\mathrm{F}}) > 1,\tag{7.116}
\end{equation*}
时非磁态不稳定——这就是 Stoner 判据。

解释为什么 $x=0$ 处 $\mathrm{d}^{2}E/\mathrm{d}x^{2}<0$ 时非磁态不稳定，并证明
这等价于上式。
